% Created 2018-09-06 Thu 16:10
% Intended LaTeX compiler: pdflatex
\documentclass[11pt]{article}
\usepackage[utf8]{inputenc}
\usepackage[T1]{fontenc}
\usepackage{graphicx}
\usepackage{grffile}
\usepackage{longtable}
\usepackage{wrapfig}
\usepackage{rotating}
\usepackage[normalem]{ulem}
\usepackage{amsmath}
\usepackage{textcomp}
\usepackage{amssymb}
\usepackage{capt-of}
\usepackage{hyperref}
\usepackage[margin=0.80in]{geometry}
\usepackage[backend=biber, style=ieee, url=false]{biblatex}
\usepackage{float}
\usepackage[super,negative]{nth}
\usepackage[capitalise]{cleveref}
\usepackage{pst-node,transparent,ragged2e}
\addbibresource{/home/nuk3/.spacemacs.d/org-files/bibliography.bib}
\DeclareFieldFormat[inproceedings]{citetitle}{\textit{#1}}
\DeclareFieldFormat[inproceedings]{title}{\textit{#1}}
\DeclareFieldFormat[misc]{citetitle}{#1}
\DeclareFieldFormat[misc]{title}{#1}
\renewcommand*{\bibpagespunct}{%
\ifentrytype{inproceedings}
{\addspace}
{\addcomma\space}}
\AtEveryCitekey{\ifuseauthor{}{\clearname{author}}}
\AtEveryBibitem{\ifuseauthor{}{\clearname{author}}}
\SpecialCoor
\author{Nyameko Lisa}
\date{31 August 2018}
\title{IPSP078 - Legal Aspects of Electronic Commerce\\\medskip
\large Assignment 1: 641343}
\hypersetup{
 pdfauthor={Nyameko Lisa},
 pdftitle={IPSP078 - Legal Aspects of Electronic Commerce},
 pdfkeywords={},
 pdfsubject={},
 pdfcreator={Emacs 27.0.50 (Org mode 9.1.2)},
 pdflang={English}}
\begin{document}

\maketitle
\addvspace{110pt}
\centering{
\pnode(0.5\textwidth,-0.5\textheight){thisCenter}
\rput(thisCenter){%\transparent{0.25}
\includegraphics[width=2.7in]{/home/nuk3/course/llb/wipo-unisa/UNISACoatofArms.eps}}}

\justifying
\addvspace{110pt}
\section*{Declaration}
\label{sec:org77e1330}
I know that plagiarism is to use someone else’s work and pass it off as my own.\\
I know that plagiarism is wrong.\\
I confirm that this assignment is my own work.\\
I have acknowledged in the bibliography accompanying the assignment all the sources that I have used.\\
I have not directly copied without acknowledgement anything from the Internet or from any other source.\\
I have indicated every quotation and citation in a footnote or bracket linked to that quotation.\\
I have not allowed anyone else to copy my work and to pass it off as their own work.\\
I understand that if any unacknowledged copying whatsoever appears in my assignment I will receive zero per cent for the assignment.\\
I am aware of the UNISA policy on plagiarism and understand that disciplinary proceedings can be instituted against me by UNISA if I contravene this policy.\\
I indicate my understanding and acceptance of this declaration by
entering my name here under:
\begin{itemize}
\item Name: \textbf{Nyameko Lisa} (Student Number: \textbf{7874-909-3})
\end{itemize}

\subsection*{NOTE}
\label{sec:orge555964}
Please note that footnotes will be denoted as \footnote{This is a footnote.} and will
appear at the bottom of the page.\\
References will be denoted by \cite{rsa02_elect_comm_trans_act} and will appear at the end of the document.
\newpage


\section{Question 1 [15]}
\label{sec:org02889b0}
\subsection{List five legal challenges that electronic commerce poses (5).}
\label{sec:org1a53eca}

Electronic commerce poses a number of regulatory and legislative challenges,
which as described in
\cite{rsa02_elect_comm_trans_act,rsa12_elect_comm_trans_amend_bill,eiselen16_ecomm,papadopoulus12_intro_cyberlaw,buys00_towards_ecomm_policy_sa,pistorius16_copyright_law_it,hamann14_direct_marketing_and_spam_via_ecomms,pistorius16_domain_names_infringe_tm}, include but are
not limited to:

\begin{itemize}
\item Contract Laws, Evidence, Validity of Agreements
\item Liability, Legal Formalities, Automated Contracts and Agency, in terms of
Time, Place, Written, Signatory and Notary Requirements of Contracts
\item Consumer Protection, Offer $\backslash$ Acceptance and Privacy Issues
\item Applicable International Guidelines, Laws and Jurisdiction
\item Harmonisation of ICT Policies Across Competitive Regional Markets
\item Fair Access to the Mechanisms and Infrastructure of Electronic Commerce
\item Intellectual Property Rights
\item Legislation of Criminal Activities Associated to Electronic Commerce
\item Taxation, Banking Laws and Electronic Payments (Cyber Credit and
Crypto-Currencies as forms of Legal Tender)
\end{itemize}

The above considerations were succinctly articulated and categorised into areas
of concern requiring serious consideration in e-commerce policy formulation in
\cite{rsa00_green_paper_ecomm_sa}, as the need for confidence in the security and
privacy of transactions performed electronically, the need to enhance the
information infrastructure for electronic commerce, the need to establish rules
that will govern electronic commerce, and the need to bring the opportunities of
e-commerce to the entire population.

\subsection{Explain the extent to which the Electronic Communication and Transactions Act of 2002 has addressed these challenges (10).}
\label{sec:orgb624b9b}

The extent to which the Act addresses some of the challenges mentioned above will be
discussed under the following headings:

\subsubsection{National e-strategy and electronics transactions policy}
\label{sec:orgf5adf67}

As per the provisions of [section 2(1)]\cite{rsa02_elect_comm_trans_act}, together
with [section 2]\cite{rsa12_elect_comm_trans_amend_bill}, the objectives of the
Act are to enable and facilitate electronic communications and transaction in
the public interest.\\

As per the provisions of [section 5(3)]\cite{rsa02_elect_comm_trans_act}, together
with [section 3]\cite{rsa12_elect_comm_trans_amend_bill}, national e-Strategy and
policy framework has been well articulated in that strict timelines have been laid out the
implementation, submission and acceptance by Parliament thereof. The Act also
permits liaison and consultation with public and private bodies in this
endeavour.\\

It is unfortunate indeed to elaborate on these severe shortcomings of the ECT
Act \cite{rsa02_elect_comm_trans_act} and the failures of Government. Initially
the Act made explicit provisions that mandated the Ministry to promote universal
access particularly to historically previously disadvantaged persons and
communities, [section 6 \& 7]\cite{rsa02_elect_comm_trans_act}. These provisions
were subsequently deleted in the amendment bill [section
4]\cite{rsa12_elect_comm_trans_amend_bill}.\\

The ECT Act \cite{rsa02_elect_comm_trans_act} has also failed to support SMMEs in
that explicit provisions outlined in the support of programmes and
infrastructure promoting the utilisation and development of electronic commerce
by SMMEs [section 9]\cite{rsa02_elect_comm_trans_act}, were subsequently removed
in the amendment bill [section 6]\cite{rsa12_elect_comm_trans_amend_bill}.

\subsubsection{Facilitation of electronic transactions}
\label{sec:org0679e14}
The ECT Act \cite{rsa02_elect_comm_trans_act} makes provisions for the legal
recognition of data messages, formation and validity of agreements and their
time and place, and the offer and acceptance of agreements, as per [sections 11,
22-24 \& 26]\cite{rsa02_elect_comm_trans_act} and [section
8]\cite{rsa12_elect_comm_trans_amend_bill}, respectively. Moreover as per
[sections 12, 13 \& 14]\cite{rsa02_elect_comm_trans_act}, the Act makes provision
for the recognition of electronic data messages as legal agreements or
contracts, in outlining the requirements data messages need satisfy in terms of
being in writing, signatures and originality, respectively.\\

As per the provisions of [sections 15-20]\cite{rsa02_elect_comm_trans_act}, together
with [section 9]\cite{rsa12_elect_comm_trans_amend_bill}, information in the form
of an electronic data message is admissible in legal proceedings and must be
given due evidential weight.

\subsubsection{Consumer protection and protection of personal information}
\label{sec:orgc3823c5}
As per the provisions of [sections 42-48 \& 50-51]\cite{rsa02_elect_comm_trans_act}
together with [sections 21-25 \& 26]\cite{rsa12_elect_comm_trans_amend_bill}, the
ECT Act \cite{rsa02_elect_comm_trans_act} has gone to great lengths to articulate
in meticulous detail both `consumer protection' and `protection of personal
information' respectively.\\

Regarding consumer protection, the provisions are more than adequate in terms of
describing: scope of applicability [section 42]\cite{rsa02_elect_comm_trans_act},
required information [section 43]\cite{rsa02_elect_comm_trans_act}, cooling-off
periods [section 44]\cite{rsa02_elect_comm_trans_act}, unsolicited goods, services
or communications [section 45]\cite{rsa02_elect_comm_trans_act}, roles and
performers [section 46]\cite{rsa02_elect_comm_trans_act}, applicability of
protection irrespective of foreign law and non-exclusion [section 47 \&
48]\cite{rsa02_elect_comm_trans_act}.\\

Regarding the protection of personal information, again the provisions are more
than adequate in terms of describing the scope of protection [section
50]\cite{rsa02_elect_comm_trans_act}, the principles for electronically collecting
personal information, together with those provisions stipulated in
\cite{rsa13_popi}.

\subsubsection{Domain name authority and administration}
\label{sec:org241f046}
As per the provisions of [chapter X]\cite{rsa02_elect_comm_trans_act}, together
with [sections 28-35]\cite{rsa12_elect_comm_trans_amend_bill}, challenges with
respect to the establishment and incorporation of a ``.za'' domain name
authority are more than adequately addressed. Moreover together with
\cite{rsa12_ect_adr}, the ECT Act \cite{rsa02_elect_comm_trans_act} makes provisions
for alternative mechanisms for the resolution of disputes arising in the ``.za''
domain name space.

\subsubsection{Limitation of liability of service providers and cyber crime}
\label{sec:org6364631}
As per the provisions of the ECT Act \cite{rsa02_elect_comm_trans_act}, a service
provider is not liable for damages arising from said service provider for
providing access to or for operating facilities as a \emph{`mere conduit'} [section
73]\cite{rsa02_elect_comm_trans_act} together with [section
38]\cite{rsa12_elect_comm_trans_amend_bill}, \emph{`caching purposes'} [section
74]\cite{rsa02_elect_comm_trans_act} together with [section
39]\cite{rsa12_elect_comm_trans_amend_bill} and \emph{`hosting services'} [section
75]\cite{rsa02_elect_comm_trans_act}. Moreover the Act explicitly excludes service
providers' obligation from monitoring data, as per [section
78]\cite{rsa02_elect_comm_trans_act}.\\

As per [chapter XII]\cite{rsa02_elect_comm_trans_act}, together with [sections
43-45]\cite{rsa12_elect_comm_trans_amend_bill}, the Act makes for adequate
provision in terms of ``cyber inspectors''. [Section
86]\cite{rsa02_elect_comm_trans_act} together with [section
48]\cite{rsa12_elect_comm_trans_amend_bill} addresses issues related to
unauthorised access to, interception of or interference with data; [section
87]\cite{rsa02_elect_comm_trans_act} together with [section
49]\cite{rsa12_elect_comm_trans_amend_bill} addresses issues pertaining to
computer related extortion, fraud and forgery; [section
88]\cite{rsa02_elect_comm_trans_act} together with [section
50]\cite{rsa12_elect_comm_trans_amend_bill} address issues pertaining to
(attempted) aiding and abetting; lastly [section
89]\cite{rsa02_elect_comm_trans_act} together with [section
51]\cite{rsa12_elect_comm_trans_amend_bill} pertains to the consequences and
penalties that a person convicted of a cyber offence is liable to.

\subsubsection{General provisions as the apply to jurisdiction of courts and saving of common law}
\label{sec:org9509a25}
As per the provisions of [section 3]\cite{rsa02_elect_comm_trans_act}, the Act
cannot be interpreted in anyway such as to exclude the application of existing
statutory or common law and must necessarily take them into account as per the
provisions of [section 3(d) \& section
10(2)(b)(iii-iv)]\cite{rsa12_elect_comm_trans_amend_bill}. Those sections also
stipulate that the Act must adhere to international best practice and conform to
the laws and guidelines of other jurisdictions and international bodies.
Finally, [chapter XIV]\cite{rsa02_elect_comm_trans_act}, together with [section
52-55]\cite{rsa12_elect_comm_trans_amend_bill} addresses issues pertaining to
jurisdiction of the courts of the Republic matters in trying offences in terms
of the ECT Act.

\section{Compare and contrast the challenges addressed in the ECT Act with the following [10]:}
\label{sec:orga75fa69}

\subsection{The UNCITRAL Model Law on Electronic Commerce of 1996 (5)}
\label{sec:orgbc6a4d5}

Challenges related to the facilitation of electronic transactions are handled in a very
similar manner in [parts 1 \& 2]\cite{rsa02_elect_comm_trans_act} \emph{`legal
requirements for data messages'} and \emph{`communication of data messages'}, and in
[chapters II and III]\cite{un96_uncitral_model_law_ecomm} \emph{`application of legal
requirements to data messages'} and \emph{`communication of data messages'}
respectively.\\

Where there are discrepancies however is in the challenges posed by matters
related to consumer protection, protection of personal information, domain name
administration, limitation of liability in respect of service providers, and
cyber crime, to name but a few. These challenges whilst addressed in the ECT Act
\cite{rsa02_elect_comm_trans_act}, are not handled at all in the UNCITRAL Model Law
\cite{un96_uncitral_model_law_ecomm}.\\

Lastly where the ECT Act \cite{rsa02_elect_comm_trans_act} has clear shortcomings,
are challenges posed by electronics commerce with regards to the carriage of
goods [articles 16 \& 17]\cite{un96_uncitral_model_law_ecomm}, which is not handled
at all within the ECT Act \cite{rsa02_elect_comm_trans_act}.

\subsection{The SADC Model Law on Electronic Transactions and Electronic Commerce of 2012. In your answer briefly explain what impact the SADC Model will have on the South African legislation and those of neighbouring countries (5)}
\label{sec:orge58a387}

The provisions of the ECT Act \cite{rsa02_elect_comm_trans_act} are almost
identically aligned with the provisions of the SADC Model Law
\cite{hipssa13_elec_trans_ecom_sadc_model_law}. From the SADC Model Law, part 1 -
\emph{`general enabling provisions'}, part 2 - \emph{`electronic transactions'}, part 3 -
\emph{`electronic commerce'}, part 4 - \emph{`consumer protection'} and part v - \emph{`service
providers'}. The SADC Model however, makes no references to domain name
authorities and cyber law to name but a few discrepancies.\\

The impact of the SADC Model, will be the harmonisation of regional ICT
communication, transaction and commerce policy. Without the SADC Model as an
overarching and guiding template, neighbouring countries may potentially develop
contradictory regulatory frameworks that impede regional electronic commerce.
\printbibliography
\end{document}